\chapter{ПРОЕКТИРОВАНИЕ И РЕАЛИЗАЦИЯ МОДУЛЬНЫХ ТЕСТОВ} \label{ch3}

% не рекомендуется использовать отдельную section <<введение>> после лета 2020 года
%\section{Введение} \label{ch3:intro}

В данной главе рассматривается проверка качества разработанных модульных тестов с помощью фреймворка мутационного тестирования Pitest.


\section{Запуск Pitest и используемые настройки} \label{ch3:sec1}

Для мутационного тестирования в проект был подключён плагин pitest-maven с поддержкой JUnit 5 (pitest-junit5-plugin).
Запуск выполнялся командой:
\texttt{mvn org.pitest:pitest-maven:mutationCoverage}

Анализ выполнялся для пакетов:
com.ipfilter.filter - модуль фильтрации
com.ipfilter.service - модуль сервисов

В результате работы Pitest формируется HTML-отчёт, содержащий сводные метрики качества тестов, покрытие строк кода, список сгенерированных мутантов, статус каждого мутанта (KILLED \/ SURVIVED \/ NO\_COVERAGE), детализацию по пакетам, классам и строкам кода.

\section{Задействованные мутационные операторы}\label{ch3:sec2}
По отчёту Pitest в проекте были использованы стандартные операторы мутаций:

\section{3.2 Задействованные мутационные операторы}

По отчёту Pitest в проекте были использованы следующие операторы мутаций:

\begin{itemize}
	\item CONDITIONALS\_BOUNDARY - Изменение граничных сравнений
	\item NEGATE\_CONDITIONALS - Отрицание условий
	\item TRUE\_RETURNS - Подмена булевых значений
	\item FALSE\_RETURNS - Подмена булевых значений
	\item NULL\_RETURNS - Подмена возвращаемого значения
	\item EMPTY\_RETURNS - Возврат пустой коллекции
	\item INCREMENTS - Изменение инкрементов
	\item MATH - Изменение математических операций 
	\item VOID\_METHOD\_CALLS - Удаление вызова void-метода
	\item CONSTRUCTOR\_CALLS - Удаление вызова конструктора
	\item INLINE\_CONSTS - Подмена встроенных констант
	\item INVERT\_NEGS - Инверсия отрицаний
	\item PRIMITIVE\_RETURNS - Подмена примитивных значений
\end{itemize}

Наиболее значимыми для данного проекта являются мутации условий и возвращаемых значений, так как они затрагивают основную логику фильтрации IP-адресов и формирования рекомендаций.

\begin{itemize}
    \item \textbf{CONDITIONALS\_BOUNDARY} - критичен для проверки граничных значений (размер фильтра, количество элементов);
    \item \textbf{NEGATE\_CONDITIONALS} - проверяет корректность условий в методах contains() и add();
    \item \textbf{NULL\_RETURNS} - тестирует обработку null-значений в методах фильтра;
    \item \textbf{VOID\_METHOD\_CALLS} - проверяет вызовы методов записи статистики и рекомендаций;
    \item \textbf{MATH} - важен для расчёта вероятности False Positive в Bloom Filter.
\end{itemize}

\noindent Данные операторы позволяют оценить качество тестов по следующим аспектам:
\begin{enumerate}
    \item Полнота проверки граничных условий;
    \item Корректность обработки исключительных ситуаций;
    \item Надёжность вычислений и расчётов;
    \item Полнота проверки возвращаемых значений.
\end{enumerate}

\section{Подробный анализ мутаций для участка кода} \label{ch3:sec3}
\subsection{Метод contains (класс BloomFilter)}
\texttt{public boolean contains(String ip)}
Применённые мутации:
\begin{itemize}
	\item NEGATE\_CONDITIONALS - инверсия if (ip == null) на if (ip != null)
	\item CONDITIONALS\_BOUNDARY - изменение i < hashFunctionsCount на i <\= hashFunctionsCount
	\item FALSE\_RETURNS - подмена return false на return true
	\item TRUE\_RETURNS - подмена return true на return false
\end{itemize}

\subsection{Метод processAndCompare (класс FilterService)}
\texttt{public OutputData processAndCompare(List<String> ipsToCheck)}
Применённые мутации:
\begin{itemize}
	\item NULL\_RETURNS — подмена результата на null
	\item EMPTY\_RETURNS — возврат пустого списка результатов
	\item VOID\_METHOD\_CALLS — удаление вызова setStatistics()
\end{itemize}

\section{Эквивалентные мутации} \label{ch3:sec4}
Эквивалентные мутации — это изменения исходного кода, которые формально изменяют текст программы, но не изменяют её фактическое поведение.
Пример из проекта:
// Оригинал
if (size > 0) {
    return calculateProbability();
}

// Мутация (CONDITIONALS\_BOUNDARY)
if (size >= 0) {
    return calculateProbability();
}

Мутация эквивалентна, потому что переменная size (размер битового массива) инициализируется положительным числом в конструкторе и никогда не бывает отрицательной. Условие size == 0 теоретически возможно только при некорректном создании фильтра, что предотвращается конструктором. Поэтому мутация не влияет на наблюдаемое поведение программы.

Pitest не определяет эквивалентные мутации автоматически, поэтому их анализ выполнялся вручную по списку выживших мутантов (SURVIVED).

\section{Метрики мутационного тестирования и их интерпретация}\label{ch3:sec5}
По итоговому отчёту Pitest для проекта получены следующие значения:
\begin{itemize}
	\item Number of Classes - 4
	\item Line Coverage - 78\%
	\item Mutation Coverage - 80\%
	\item Test Strength - 93\%
\end{itemize}

Полученные значения показывают, что:
– тестами покрыта большая часть строк кода (78\%);
– тесты обнаруживают большую часть сгенерированных мутаций (80\%);
– качество тестов на покрытом коде остаётся высоким (93\%).

\section{Оценка mutation-adequate и улучшение качества тестов}\label{ch3:sec6}
С учётом полученных метрик можно сделать вывод, что разработанный набор модульных тестов имеет хорошее качество и в целом удовлетворяет требованиям мутационной адекватности (mutation-adequate) для учебного проекта.
Рекомендации по улучшению:
\begin{itemize}
	\item Добавить тесты на граничные значения (пустой список, 1 элемент, максимум)
	\item Добавить тесты на обработку null и пустых строк
	\item Добавить тесты на все ветви if/else в методах фильтрации
	\item Добавить тесты на расчёт статистики и вероятности False Positive
	\item Добавить тесты на метод generateRecommendation() с разными сценариями
\end{itemize}